% !TeX root = ../../Book2.tex
% 有限先修与课程选择集中说明，不恢复纯导航表。
\ProseParagraphBreak
\noindent\textbf{第一卷的有限回查点。}
基本积分与 \(L^p\) 理论贯穿本卷，较专门的输入只在相应位置需要。
\begin{itemize}
\item 进入\chapref{b2:ch:10}前，补读第一卷定理 8.5.5 的紧自伴谱定理，以及命题 21.4.1 的区间 Sobolev 绝对连续代表。前者在第一卷属于选读，在本章却是 Green 逆算子谱展开的必需工具。本章先从常微分方程构造紧逆，不借用后面的 Lax--Milgram 定理；未准备好谱论时，可以先读全空间波动方程与热核，暂缓区间谱展开。
\item \chapref{b2:ch:14}使用第一卷定理 9.4.7 的弱收敛子列、定理 9.5.2 的闭凸集弱闭性，以及定理 9.5.8--9.5.9 的凸泛函弱下半连续性与直接方法。用于梯度能量时还须核对 \(1<p<\infty\) 的自反性和零边界 Poincaré 不等式，不能只记住“有界序列可取子列”而忘记空间及收敛方式。
\item \chapref{b2:ch:15}使用第一卷定理 8.4.7 的 Hilbert 空间 Riesz 表示，并在 \(H_0^1(\Omega)\) 中求弱解。第一卷定义 22.2.1 以光滑紧支集函数的闭包规定零边界条件；在有界 Lipschitz 区域上，才按第一卷定理 23.1.6 解释为零迹。强制性所需的 Poincaré 不等式及有关嵌入见第一卷第二十四章。
\item \chapref{b2:ch:16}沿用上述能量空间，按所用估计回查 Sobolev 不等式、截断与局部坐标下的弱导数。内部估计不要求整个边界光滑；全局二阶 Dirichlet 估计则使用正文列明的 \(C^{1,1}\) 边界和 Lipschitz 系数，两者不能混同。
\item 第十五章末的 Fredholm 选读部分另用第一卷附录 C 的紧扰动理论及推论 25.2.6 的零边界紧嵌入。\appref{b2:app:D}的一般非负自伴算子视角还需要无界自伴谱定理与 Borel 函数演算，其范围大于第一卷的紧自伴谱定理。附录另给准确来源；本卷热方程的 Fourier 证明不依赖这一抽象视角。
\end{itemize}

\ProseParagraphBreak
\noindent\textbf{本卷所需的复分析工具。}
\secref{b2:sec:0402}在本卷内证明 Fourier 插值所需的特殊三线估计，读到这里不必先学完整复分析。\secref{b2:sec:1704}处理一般解析插值时，使用全纯函数的局部光滑性、Cauchy--Riemann 方程及实虚部调和，并回查\cref{b2:cor:13-boundary-maximum}的调和最大值原理；沿调和分析路线跳过第二编时，只需补读\secref{b2:sec:1301}中这项结果，不必先读全部边值问题。

\ProseParagraphBreak
选读的\chapref{b2:ch:20}另用局部 Cauchy 公式、唯一性和 Cauchy 导数估计。即使跳过该章，进入普通正文\secref{b2:sec:2304}前也须掌握局部 Cauchy 公式及其局部一致极限、幂级数和唯一性推论；该节会回忆这些事实，再构造所需的复 Gamma 函数和分布解析族。这里把局部 Cauchy 理论作为基础复分析先修，不要求先学留数计算、模形式或整套第三卷。复参数出现在公式中，并不意味着参数求导与取极限已经合法，正文仍逐处核验控制条件。

\ProseParagraphBreak
\noindent\textbf{课程可以怎样取舍？}
\begin{itemize}
\item Fourier 基础课程以第一编的级数、反演、Plancherel、插值和分布为主体。第五章的 Hermite 专题及第六章相关补例可先略过，进入\appref{b2:app:B}前补读；Poisson 求和、采样、周期热核和 theta 的算术联系可按用途选择。
\item 经典与弱解 PDE 课程从第一编所需工具进入第七至十三章，再读第十四至十六章。区间谱理论按前述方式准备；\appref{b2:app:D}接热方程，\appref{b2:app:H}接椭圆弱解。直接法、非对称双线性型和正则性分别解决不同问题，不能以其中一个替代另外两个。
\item 实变调和分析课程从第一编进入第十七至十九章，再读第二十二章，并按上文补齐调和最大值原理。非齐次分解及 \(F^s_{p,2}\) 的最小定义放在\secref{b2:sec:1904}，一般 Besov、Triebel--Lizorkin 选读节可以暂缓。第二十二章的基本乘子证明不依赖驻相，可以先于第二十一章学习；第二十章是独立选读支线。
\item 振荡与限制课程在上述调和工具之后读第二十一章，按需补\appref{b2:app:A}的径向 Fourier、Bessel 渐近和幂权核公式，再读第二十三、二十四章，最后选读第二十五章。第二十三章的解析构造须先补足前述局部复分析；\appref{b2:app:E}的色散方程与\appref{b2:app:G}的微局部方法是进一步的不同方向。
\end{itemize}

\ProseParagraphBreak
\noindent\textbf{选读与证明深度分开判断。}
编号右侧的星号表示可以从相应普通阅读路线中略去；附录整体也是选读材料，但更深入的路线可能需要其中的特定引理，使用处会明确回指。星号不表示证明简略：第二十章有完整的模型证明，普通正文也可能引用注明来源的深层定理。应分别辨认完整证明、说明缺失步骤的证明概要，以及有明确假设和出处的外部输入，不以“普通正文”或“选读”代替证明责任的判断。
